\appendix
\renewcommand{\thechapter}{\Asbuk{chapter}}
\titleformat{\chapter}[block]{\normalfont\large\bfseries\centering}{Приложение \thechapter.}{0.5em}{}
\addtocontents{toc}{\protect\renewcommand{\protect\cftchappresnum}{Приложение }}
\addtocontents{toc}{\protect\setlength{\protect\cftchapnumwidth}{8em}}

\chapter{Воспроизведение экспериментов}

Исходный код, конфигурация стенда, сценарии экспериментов и сырые данные
всех прогонов находятся в репозитории \url{https://github.com/hanglider/cursach}.
Требования к машине: docker с виртуальной машиной на 10~ГБ памяти.

\begin{table}[ht]
  \centering
  \caption{Основные действия}
  \small
  \begin{tabular}{@{}lll@{}}
    \toprule
    Действие & Команда & Время \\
    \midrule
    Развернуть стенд & \texttt{make stand-up} & около 3 мин \\
    Серия сценария 1 & \texttt{make dr-series} & около 1.5 ч \\
    Серия сценария 2 & \texttt{make sched-series} & 1.5--2 ч \\
    Серия сценария 3 & \texttt{make policy-series} & около 35 мин \\
    Анализ и графики & \texttt{make analysis}, \texttt{make plots} & минуты \\
    \bottomrule
  \end{tabular}
\end{table}

\chapter{Структура репозитория}

\begin{table}[ht]
  \centering
  \caption{Каталоги репозитория}
  \small
  \begin{tabularx}{\textwidth}{@{}lX@{}}
    \toprule
    Каталог & Содержимое \\
    \midrule
    \texttt{stand/} & Развёртывание стенда: кластеры, наблюдаемость, шлюз,
      политики \\
    \texttt{test-service/} & Тестовый сервис, Go \\
    \texttt{registry/} & Каталог кластеров, Python, FastAPI, PostgreSQL \\
    \texttt{scheduler/} & Планировщик уровня федерации и наивный
      планировщик, Go \\
    \texttt{dr-controller/} & Контроллер аварийного восстановления, Go \\
    \texttt{experiments/} & Сценарии отказов и нагрузки, сырые результаты \\
    \texttt{analysis/} & Обработка результатов, графики, проверка гипотез \\
    \texttt{docs/} & Постановка задачи, архитектура, принятые решения \\
    \bottomrule
  \end{tabularx}
\end{table}
